% ============================================================
%  Paper 29: Modified Photon Dispersion and Cosmic Gamma-Ray Delays in SDGFT
%  Target: RevTeX 4-2 Compilation (SDGFT Series)
% ============================================================
\documentclass[amsmath,amssymb,aps,prd,twocolumn,superscriptaddress,nofootinbib]{revtex4-2}

\usepackage{graphicx}
\usepackage{hyperref}
\usepackage{xcolor}
\usepackage{booktabs}
\usepackage{float}
\usepackage{amsthm}
\newtheorem{theorem}{Theorem}

% ── SDGFT shared macros ──────────────────────────────────────
\usepackage{sdgft-macros}

% ── Document ─────────────────────────────────────────────────
\begin{document}

\preprint{SDGFT-2026-29}

\title{Modified Photon Dispersion and Cosmic Gamma-Ray Delays in SDGFT}

\author{David A. Besemer}
\email{david.besemer@sdgft.org}
\affiliation{Independent Researcher}

\date{\today}

\begin{abstract}
We derive the modified photon dispersion relation from the discrete 24-cell lattice structure of spacetime. We calculate the cosmological propagation delays of high-energy photons, showing that the LIV parameter is xi_LIV = sqrt(24/67) ~ 0.5985. We compare our results to GRB 090510, showing that the predicted delay of 0.63 s complies with Fermi-LAT limits.
\end{abstract}

\maketitle

\section{Introduction}
Lorentz invariance violation (LIV) is a common prediction of quantum gravity. We derive the modified dispersion relation for photons in SDGFT.

\section{Theoretical Formulation}
Here we formulate the core mathematical relations governing the system:
\begin{equation}
  E^2 = p^2 c^2 \left(1 + \xi_{\text{LIV}} \frac{E}{\EPl}\right) \quad (\text{modified dispersion relation})
\end{equation}
\begin{equation}
  \Delta t = \xi_{\text{LIV}} \frac{E}{\EPl} \int_0^z \frac{1+z'}{H(z')} dz' \approx 0.63 \text{ s} \quad (\text{for GRB 090510 at 31 GeV})
\end{equation}
\section{Gamma-Ray Burst Delay}
We integrate the propagation delay over cosmological distances under the Horndeski background, showing that the delay is within the Fermi-LAT upper limits.

\section{Conclusion}
The predicted delay is consistent with Fermi-LAT and MAGIC observations, and will be tested by the Cherenkov Telescope Array (CTA).



\begin{thebibliography}{99}
\bibitem{Goldstein_Paper01} D. A. Besemer, ``Axiomatic Foundations of SDGFT,'' SDGFT-2026-01 (in preparation).
\bibitem{Goldstein_Paper04} D. A. Besemer, ``Effective Spectral Dimension and the Banach Fixed-Point Theorem in SDGFT,'' SDGFT-2026-04.
\bibitem{Goldstein_Code} D. A. Besemer, \texttt{sdgft} Python package, \url{https://github.com/david-besemer/sdgft} (2026).
\end{thebibliography}

\end{document}
